\paragrafos{Συμμετρία ως προς άξονα: ο καθρέφτης του κήπου}{Β.2.1, «Συμμετρία ως προς άξονα»}

\begin{anagke}
Η Αγγελική σχεδιάζει πια το δεύτερο σχέδιο του κήπου: έναν κήπο με κεντρικό ευθύγραμμο μονοπάτι, και δύο όμοια παρτέρια, ένα δεξιά και ένα αριστερά του. «Θέλω να είναι το ένα η εικόνα του άλλου στον καθρέφτη», λέει. Ο Στέλιος ρωτάει: «Ποιο σημείο του αριστερού παρτεριού αντιστοιχεί στο σημείο του δεξιού;». Η Αγγελική δοκιμάζει ένα κόλπο: σχεδιάζει το αριστερό παρτέρι στο φύλλο, διπλώνει το φύλλο κατά μήκος του μονοπατιού και με μια σταθερή πίεση αντιγράφει το σχήμα στην άλλη πλευρά.

«Κάθε σημείο πέφτει ακριβώς πάνω σε ένα σημείο της άλλης πλευράς!», λέει ο Νίκος. Η Μυρτώ το ελέγχει με τον χάρακα: το δέντρο που απέχει $1{,}2$ μέτρα από το μονοπάτι έχει «ζευγάρι» ένα δέντρο που απέχει επίσης $1{,}2$ μέτρα, από την άλλη πλευρά και «ακριβώς απέναντι», πάνω στην κάθετη. Η Πολυξένη γράφει στον πίνακα: «Συμμετρία».

Σε αυτή την παράγραφο ορίζουμε ποιο είναι το συμμετρικό ενός σημείου και ενός σχήματος ως προς μια ευθεία, και ανακαλύπτουμε ότι η συμμετρία διατηρεί τα μήκη και τις γωνίες.
\end{anagke}

\begin{lexi}{συμμετρία, άξονας}{συμμετρία, άξονας}
Η λέξη \lexiword{συμμετρία} (σύν + μέτρο) σημαίνει «κοινό μέτρο»: ό,τι έχει ίσα μέτρα από δύο μεριές. Ο \lexiword{άξονας} είναι η ευθεία γύρω από την οποία «στρέφεται» ο καθρέφτης, αυτός που χωρίζει τα δύο συμμετρικά μέρη.
\end{lexi}

\begin{metafrasi}
Το σημείο που πέφτει πάνω στο $\mathrm{Β}$ όταν διπλώνω. & το συμμετρικό $\mathrm{Β}'$ του $\mathrm{Β}$ ως προς την $\varepsilon$\\
Ίση απόσταση, κάθετα. & $\mathrm{ΒΒ}'\perp\varepsilon$ και η $\varepsilon$ περνά από το μέσο του $\mathrm{ΒΒ}'$\\
Σημείο πάνω στον άξονα. & είναι συμμετρικό του εαυτού του\\
Ο καθρέφτης δεν αλλάζει μήκη. & $\mathrm{ΑΒ}=\mathrm{Α}'\mathrm{Β}'$\\
\end{metafrasi}

\tmima{Η θεωρία}

\begin{orismos}{συμμετρικό σημείου ως προς ευθεία}
\textbf{Συμμετρικό} του σημείου $\mathrm{Β}$ ως προς την ευθεία $\varepsilon$ είναι το σημείο $\mathrm{Β}'$ με το οποίο συμπίπτει το $\mathrm{Β}$ όταν διπλώσουμε το επίπεδο κατά μήκος της $\varepsilon$. Η $\varepsilon$ λέγεται \textbf{άξονας συμμετρίας}. Κάθε σημείο της $\varepsilon$ είναι συμμετρικό του εαυτού του.
\end{orismos}

Από τη δίπλωση προκύπτει ότι το τμήμα $\mathrm{ΒΒ}'$ είναι \textbf{κάθετο} στην $\varepsilon$ και η $\varepsilon$ περνά από το \textbf{μέσο} του: τα $\mathrm{Β}$ και $\mathrm{Β}'$ απέχουν το ίδιο από την $\varepsilon$, από δύο αντίθετες μεριές. Έτσι κατασκευάζουμε το συμμετρικό με \emph{γνώμονα}: φέρνουμε την κάθετη από το $\mathrm{Β}$ προς την $\varepsilon$, και στην προέκτασή της παίρνουμε σημείο $\mathrm{Β}'$ με την ίδια απόσταση από την $\varepsilon$.

\begin{center}
\begin{tikzpicture}[scale=1.1]
\draw[thick,{Stealth}-{Stealth}] (0,-0.2)--(0,3.4) node[above] {$\varepsilon$};
\fill (-1.6,2.5) circle (2pt) node[left] {$\mathrm{Β}$}; \fill (1.6,2.5) circle (2pt) node[right] {$\mathrm{Β}'$};
\draw[dashed,gblue] (-1.6,2.5)--(1.6,2.5);
\fill (-1.2,0.6) circle (2pt) node[left] {$\mathrm{Α}$}; \fill (1.2,0.6) circle (2pt) node[right] {$\mathrm{Α}'$};
\draw[dashed,gblue] (-1.2,0.6)--(1.2,0.6);
\draw (-1.6,2.5)--(-1.2,0.6); \draw (1.6,2.5)--(1.2,0.6);
\draw (0,2.5)++(0.15,0)--++(0,-0.15)--++(-0.15,0);
\end{tikzpicture}
\end{center}

\begin{kanonas}{συμμετρικά σχήματα}
Δύο σχήματα $(\Sigma_1)$ και $(\Sigma_2)$ λέγονται \textbf{συμμετρικά ως προς την $\varepsilon$} όταν το καθένα αποτελείται από τα συμμετρικά των σημείων του άλλου. Τα συμμετρικά ως προς ευθεία σχήματα είναι \textbf{ίσα}: έχουν ίσα μήκη και ίσες γωνίες (με δίπλωση συμπίπτουν).
\end{kanonas}

Έτσι το συμμετρικό ενός ευθυγράμμου τμήματος είναι ευθύγραμμο τμήμα με το ίδιο μήκος και το συμμετρικό μιας γωνίας είναι γωνία με το ίδιο μέτρο. Για να σχεδιάσω το συμμετρικό ενός πολυγώνου, βρίσκω τα συμμετρικά των κορυφών και τα ενώνω με την ίδια σειρά.

\begin{prosoxi}
\begin{itemize}[leftmargin=1.2em,itemsep=0.3ex,topsep=2pt]
\item Το τμήμα $\mathrm{ΒΒ}'$ είναι \emph{κάθετο} στον άξονα (όχι παράλληλο).
\item Η απόσταση του $\mathrm{Β}$ από τον άξονα είναι ίση με της $\mathrm{Β}'$, όχι ίση με την απόσταση $\mathrm{ΒΒ}'$: το $\mathrm{ΒΒ}'$ είναι διπλάσιο.
\item Το συμμετρικό ενός σχήματος είναι ίσο με αυτό, αλλά «ανάποδο» (σαν καθρέφτης): τη φορά δεν την κρατάει.
\item Το συμμετρικό του $\mathrm{Β}'$ είναι πάλι το $\mathrm{Β}$.
\end{itemize}
\end{prosoxi}

\tmima{Λυμένα παραδείγματα}

\paradeigma{Το συμμετρικό σημείου}
\tixero Το δέντρο $\mathrm{Β}$ απέχει $1{,}2$ μ. από την ευθεία $\varepsilon$.\par
\tithelo Πού βρίσκεται το συμμετρικό του και πόσο απέχουν τα δύο.\par
\possindeo Το $\mathrm{Β}'$ είναι στην κάθετη από το $\mathrm{Β}$ προς την $\varepsilon$, από την άλλη μεριά, με απόσταση $1{,}2$ μ. Η απόσταση $\mathrm{ΒΒ}'=2{,}4$ μ.\par
\apantisi $\mathrm{Β}'$ απέχει $1{,}2$ μ. από την $\varepsilon$ και $\mathrm{ΒΒ}'=2{,}4$ μ.\par
\elegxos Η $\varepsilon$ είναι στο μέσο του $\mathrm{ΒΒ}'$: $1{,}2+1{,}2=2{,}4$.

\paradeigma{Σημείο πάνω στον άξονα}
\tixero Το σημείο $\mathrm{Μ}$ βρίσκεται πάνω στην $\varepsilon$.\par
\tithelo Το συμμετρικό του.\par
\possindeo Αφού διπλώνω κατά μήκος της $\varepsilon$, το $\mathrm{Μ}$ δεν κινείται.\par
\apantisi Το συμμετρικό του είναι το ίδιο το $\mathrm{Μ}$.\par
\elegxos Η απόσταση του $\mathrm{Μ}$ από την $\varepsilon$ είναι $0$.

\paradeigma{Συμμετρικό τριγώνου}
\tixero Τρίγωνο $\mathrm{ΑΒΓ}$ με $\mathrm{ΑΒ}=5$ cm, $\widehat{\mathrm{Β}}=40^\circ$ και το συμμετρικό του $\mathrm{Α}'\mathrm{Β}'\Gamma'$.\par
\tithelo Το $\mathrm{Α}'\mathrm{Β}'$ και η $\widehat{\mathrm{Β}'}$.\par
\possindeo Η συμμετρία διατηρεί μήκη και γωνίες.\par
\apantisi $\mathrm{Α}'\mathrm{Β}'=5$ cm και $\widehat{\mathrm{Β}'}=40^\circ$.\par
\elegxos Τα συμμετρικά σχήματα είναι ίσα.

\paradeigma{Κατασκευή συμμετρικού}
\tixero Η ευθεία $\varepsilon$ και το σημείο $\mathrm{Α}$ εκτός αυτής.\par
\tithelo Να κατασκευάσω το συμμετρικό του.\par
\possindeo (1) Φέρνω από το $\mathrm{Α}$ την κάθετη προς την $\varepsilon$ με γνώμονα και χάρακα. (2) Μετράω την απόσταση $\mathrm{ΑΜ}$ ως το ίχνος $\mathrm{Μ}$. (3) Στην προέκταση βρίσκω το $\mathrm{Α}'$ με $\mathrm{ΜΑ}'=\mathrm{ΜΑ}$.\par
\apantisi Το σημείο $\mathrm{Α}'$.\par
\elegxos Διπλώνω το χαρτί κατά μήκος της $\varepsilon$: το $\mathrm{Α}$ πέφτει πάνω στο $\mathrm{Α}'$.

\tmima{Ασκήσεις}

\zoni{Το κατάλαβα}{συμμετρικό σημείου και σχήματος}

\ex Σχεδίασε ευθεία $\varepsilon$ και σημείο $\mathrm{Α}$ εκτός της. Βρες το συμμετρικό του με γνώμονα και με δίπλωση.

\ex Ποιο είναι το συμμετρικό ενός σημείου που ανήκει στην $\varepsilon$;

\ex Το $\mathrm{Α}$ απέχει $3{,}5$ cm από την $\varepsilon$. Πόσο απέχει το συμμετρικό του από την $\varepsilon$ και πόσο από το $\mathrm{Α}$;

\ex Σχεδίασε ένα τρίγωνο $\mathrm{ΑΒΓ}$ και το συμμετρικό του ως προς μια ευθεία $\varepsilon$ που δεν το τέμνει.

\ex Αν στο τρίγωνο $\mathrm{ΑΒΓ}$ η γωνία $\widehat{\mathrm{Β}}$ είναι $40^\circ$ και $\mathrm{ΑΒ}=5$ cm, τι ξέρεις για το συμμετρικό τρίγωνο;

\ex Σχεδίασε το συμμετρικό ενός ευθυγράμμου τμήματος $\mathrm{ΑΒ}=4{,}2$ cm ως προς ευθεία $\varepsilon$. Πόσο είναι το μήκος του;

\ex Σωστό ή λάθος; α) Κάθε σημείο της $\varepsilon$ είναι συμμετρικό του εαυτού του. \quad β) Το $\mathrm{ΑΑ}'$ είναι παράλληλο προς την $\varepsilon$. \quad γ) Συμμετρικά σχήματα είναι ίσα. \quad δ) Το συμμετρικό του $\mathrm{Α}'$ είναι το $\mathrm{Α}$. \quad ε) Η συμμετρία αλλάζει τα μήκη.

\zoni{Το χρησιμοποιώ}{κήπος, καθρέφτες, γράμματα}

\ex \textbf{Ο καθρέφτης του κήπου.} Ένα δέντρο απέχει $1{,}2$ μ. από το κεντρικό μονοπάτι. Πού πρέπει να φυτευτεί το συμμετρικό του; Πόσο απέχουν τα δύο δέντρα;

\ex \textbf{Η λίμνη.} Η αντανάκλαση ενός δέντρου στη λίμνη είναι το συμμετρικό του ως προς την ακτή. Αν το δέντρο είναι $4$ μ. ψηλό, πόσο «βαθιά» φαίνεται η αντανάκλαση;

\ex \textbf{Το ρολόι.} Στον καθρέφτη ένα ρολόι δείχνει $3{:}00$. Τι ώρα είναι πραγματικά; (Σκέψου πού είναι ο ωροδείκτης.)

\ex \textbf{Τα γράμματα.} Ποια από τα κεφαλαία γράμματα $\mathrm{Α}$, $\mathrm{Β}$, $\mathrm{Η}$, $\mathrm{Μ}$, $\mathrm{Ο}$, $\mathrm{Τ}$, $\mathrm{Ζ}$ μένουν ίδια όταν τα καθρεφτίσεις ως προς κατακόρυφη ευθεία;

\ex \textbf{Η πεταλούδα.} Σχεδίασε στο τετράδιο μια ευθεία και το μισό μιας πεταλούδας στη μία πλευρά. Συμπλήρωσε το άλλο μισό ως συμμετρικό.

\zoni{Το σκέφτομαι}{γιατί κάθετα;}

\ex Γιατί το τμήμα $\mathrm{ΑΑ}'$ είναι κάθετο στον άξονα; (Σκέψου τη δίπλωση.)

\ex Ποια σημεία του επιπέδου μένουν στη θέση τους όταν καθρεφτίζουμε ως προς την $\varepsilon$;

\ex Αν το $\mathrm{Α}'$ είναι συμμετρικό του $\mathrm{Α}$ ως προς την $\varepsilon$, και το $\mathrm{Α}''$ συμμετρικό του $\mathrm{Α}'$, πού βρίσκεται το $\mathrm{Α}''$;

\ex Το τμήμα $\mathrm{ΑΒ}$ τέμνει την $\varepsilon$ στο $\mathrm{Μ}$. Εξήγησε γιατί το συμμετρικό $\mathrm{Α}'\mathrm{Β}'$ περνά επίσης από το $\mathrm{Μ}$.

\zoni{Για το σπίτι}{πείραμα με μελάνι}

\ex Βάλε λίγο μελάνι στη μέση ενός χαρτιού και δίπλωσέ το. Άνοιξέ το και περιέγραψε τι σχήμα βγαίνει και πού είναι ο άξονας συμμετρίας.

\ex Βρες στο σπίτι πέντε αντικείμενα που φαίνονται συμμετρικά.
